% --- CORRECCION: Forzar interpretacion raw de la entrada ---
\UseRawInputEncoding
\documentclass[aspectratio=169, 10pt, xcolor=table]{beamer}

% --- Configuracion del Tema ---
\usetheme{Madrid}
\usecolortheme{dolphin}



% --- Paquetes necesarios ---
\usepackage[utf8]{inputenc}
\usepackage[T1]{fontenc}
\usepackage[spanish,es-noquoting]{babel}
\usepackage{amsmath, amssymb}
\usepackage{mathtools}
\usepackage{booktabs}
\usepackage{graphicx}

\usepackage{listings}
\usepackage{tikz}
\usetikzlibrary{arrows.meta, positioning, shapes.geometric, babel}
\usepackage{etoolbox}
\usepackage{upquote}



% --- Estilo para codigo SQL y Python ---
\definecolor{codegreen}{rgb}{0,0.6,0}
\definecolor{codegray}{rgb}{0.5,0.5,0.5}
\definecolor{codepurple}{rgb}{0.58,0,0.82}
\definecolor{backcolour}{rgb}{0.95,0.95,0.92}
\lstdefinestyle{mysql}{
    language=SQL,
    backgroundcolor=\color{backcolour},
    commentstyle=\color{codegreen},
    keywordstyle=\color{blue},
    numberstyle=\tiny\color{codegray},
    stringstyle=\color{codepurple},
    basicstyle=\ttfamily\scriptsize,
    breaklines=true,
    frame=single,
    showstringspaces=false,
    literate={á}{{\'a}}1 {é}{{\'e}}1 {í}{{\'i}}1 {ó}{{\'o}}1 {ú}{{\'u}}1 {ñ}{{\~n}}1 {ü}{{\"u}}1
}
\lstdefinestyle{python}{
    language=Python,
    backgroundcolor=\color{backcolour},
    commentstyle=\color{codegreen},
    keywordstyle=\color{blue},
    numberstyle=\tiny\color{codegray},
    stringstyle=\color{codepurple},
    basicstyle=\ttfamily\scriptsize,
    breaklines=true,
    frame=single,
    showstringspaces=false,
    literate={á}{{\'a}}1 {é}{{\'e}}1 {í}{{\'i}}1 {ó}{{\'o}}1 {ú}{{\'u}}1 {ñ}{{\~n}}1 {ü}{{\"u}}1
}
\lstset{style=python} % Por defecto Python, aunque se puede cambiar en cada bloque

% --- Pie de pagina personalizado ---
\makeatletter
\setbeamertemplate{footline}{
  \leavevmode%
  \hbox{%
  \begin{beamercolorbox}[wd=.333333\paperwidth,ht=2.25ex,dp=1ex,center]{author in head/foot}%
    \usebeamerfont{author in head/foot}\insertshortauthor
  \end{beamercolorbox}%
  \begin{beamercolorbox}[wd=.333333\paperwidth,ht=2.25ex,dp=1ex,center]{title in head/foot}%
    \usebeamerfont{title in head/foot}Sesión 11
  \end{beamercolorbox}%
  \begin{beamercolorbox}[wd=.333333\paperwidth,ht=2.25ex,dp=1ex,right]{date in head/foot}%
    \usebeamerfont{date in head/foot}BBDD \hfill \insertframenumber/\inserttotalframenumber \hspace*{0.3cm}
  \end{beamercolorbox}}%
  \vskip0pt%
}
\makeatother

% --- Datos de la portada ---
\title[SESIÓN 11 BBDD]{\textbf{Sesión 11}}
\subtitle{NoSQL y Cloud}
\author[Prof. C. Sánchez C.]{Carlos César Sánchez Coronel}
\institute{\textbf{BASES DE DATOS}}
\date{2026}

\setbeamertemplate{navigation symbols}{}

\AtBeginSection[]{
  \begin{frame}
    \frametitle{Contenido de la sección}
    \tableofcontents[currentsection]
  \end{frame}
}

\begin{document}

\begin{frame}
    \titlepage
\end{frame}

\begin{frame}{Contenido general}
    \tableofcontents
\end{frame}

% ============================================================================
% SECCIÓN 1: INTRODUCCIÓN A NoSQL
% ============================================================================
\section{Introducción a NoSQL}

\begin{frame}{¿Por qué NoSQL?}
    \begin{itemize}
        \item Escalabilidad vertical costosa (más CPU/RAM en un solo servidor).
        \item Esquema rígido (schema-on-write) dificulta cambios frecuentes.
        \item Dificultad para manejar grandes volúmenes de datos no estructurados.
        \item Rendimiento limitado en operaciones de lectura/escritura masivas.
    \end{itemize}
\end{frame}

\begin{frame}{Teorema CAP}
    \begin{block}{En sistemas distribuidos no se pueden garantizar simultáneamente:}
        \begin{itemize}
            \item \textbf{C}onsistencia: todos ven los mismos datos.
            \item \textbf{D}isponibilidad: siempre responde.
            \item \textbf{T}olerancia a Particiones: funciona pese a fallos de red.
        \end{itemize}
    \end{block}
    \begin{description}
        \item[CP] HBase, MongoDB (configuración típica)
        \item[AP] Cassandra, DynamoDB, CouchDB
        \item[CA] Sistemas tradicionales (no toleran particiones)
    \end{description}
\end{frame}

\begin{frame}{Tipos de Bases de Datos NoSQL}
    \begin{description}
        \item[Documentales] MongoDB, Couchbase
        \item[Clave-Valor] Redis, DynamoDB
        \item[Grafos] Neo4j, Amazon Neptune
        \item[Columnares] Cassandra, HBase
        \item[Motores de búsqueda] Elasticsearch, Solr
    \end{description}
\end{frame}

% ============================================================================
% SECCIÓN 2: BASES DE DATOS DOCUMENTALES: MONGODB
% ============================================================================
\section{MongoDB}

% --- CORREGIDO: Añadido [fragile] a este frame ---
\begin{frame}[fragile]{Modelo de Datos en MongoDB}
    \begin{itemize}
        \item Almacena documentos en formato BSON (JSON binario).
        \item Esquema flexible: diferentes documentos pueden tener distintos campos.
        \item Soporta arrays, subdocumentos y tipos de datos avanzados.
    \end{itemize}
    \begin{exampleblock}{Ejemplo de documento}
        \begin{verbatim}
{
  "_id": ObjectId("507f1f77bcf86cd799439011"),
  "nombre": "Ana Pérez",
  "email": "ana@mail.com",
  "direccion": {
    "calle": "Av. Siempre Viva",
    "ciudad": "Madrid"
  },
  "pedidos": [
    { "fecha": "2025-01-10", "total": 150 },
    { "fecha": "2025-02-15", "total": 200 }
  ]
}
        \end{verbatim}
    \end{exampleblock}
\end{frame}

\begin{frame}[fragile]{Consultas con PyMongo}
    \begin{block}{Operaciones básicas}
        \begin{lstlisting}[style=python]
from pymongo import MongoClient

client = MongoClient('localhost', 27017)
db = client['tienda']
coleccion = db['clientes']

# Insertar
coleccion.insert_one({"nombre": "Luis", "email": "luis@mail.com"})

# Buscar
result = coleccion.find({"ciudad": "Madrid"})
for doc in result:
    print(doc)

# Actualizar
coleccion.update_one({"nombre": "Luis"}, {"$set": {"ciudad": "Barcelona"}})

# Eliminar
coleccion.delete_one({"nombre": "Luis"})
        \end{lstlisting}
    \end{block}
\end{frame}

\begin{frame}{Índices y Escalabilidad}
    \begin{itemize}
        \item \textbf{Índices}: simples, compuestos, geoespaciales, texto.
        \item \textbf{Sharding}: particionamiento horizontal por clave de shard.
        \item \textbf{Réplicas}: alta disponibilidad (replica sets).
    \end{itemize}
\end{frame}

% ============================================================================
% SECCIÓN 3: BASES DE DATOS CLAVE-VALOR: REDIS
% ============================================================================
\section{Redis}

\begin{frame}{Características de Redis}
    \begin{itemize}
        \item Base de datos en memoria (in-memory).
        \item Latencia de microsegundos.
        \item Estructuras ricas: strings, hashes, listas, sets, sorted sets, etc.
        \item Operaciones atómicas (INCR, LPUSH, ZADD).
    \end{itemize}
\end{frame}

\begin{frame}[fragile]{Conexión desde Python (redis-py)}
    \begin{lstlisting}[style=python]
import redis

r = redis.Redis(host='localhost', port=6379, decode_responses=True)

# Strings
r.set('clave', 'valor')
print(r.get('clave'))

# Contadores
r.incr('contador')

# Listas
r.lpush('lista', 'elemento1', 'elemento2')
elementos = r.lrange('lista', 0, -1)

# Sorted sets (para ranking)
r.zadd('ranking', {'A': 1, 'B': 2})
ranking = r.zrevrange('ranking', 0, -1, withscores=True)
    \end{lstlisting}
\end{frame}

\begin{frame}{Persistencia y Escalabilidad}
    \begin{itemize}
        \item \textbf{Persistencia}: RDB (snapshots) y AOF (log de operaciones).
        \item \textbf{Replicación}: maestro-esclavo.
        \item \textbf{Redis Cluster}: sharding automático.
    \end{itemize}
\end{frame}

% ============================================================================
% SECCIÓN 4: BASES DE DATOS DE GRAFOS: NEO4J
% ============================================================================
\section{Neo4j}

\begin{frame}{Modelo de Grafos}
    \begin{itemize}
        \item \textbf{Nodos}: entidades (Persona, Empresa).
        \item \textbf{Relaciones}: conexiones dirigidas con tipo y propiedades.
        \item \textbf{Propiedades}: atributos de nodos y relaciones.
    \end{itemize}
\end{frame}

\begin{frame}[fragile]{Consultas con Cypher}
    \begin{block}{Crear nodos y relaciones}
        \begin{lstlisting}
CREATE (a:Persona {nombre: 'Ana', edad: 30})
CREATE (b:Persona {nombre: 'Luis', edad: 32})
CREATE (a)-[:AMIGO_DE]->(b)
        \end{lstlisting}
    \end{block}
    \begin{block}{Consultar amigos de Ana}
        \begin{lstlisting}
MATCH (p:Persona {nombre: 'Ana'})-[:AMIGO_DE]->(amigo)
RETURN amigo.nombre
        \end{lstlisting}
    \end{block}
\end{frame}

\begin{frame}[fragile]{Integración con Python (neo4j driver)}
    \begin{lstlisting}[style=python]
from neo4j import GraphDatabase

driver = GraphDatabase.driver("bolt://localhost:7687", auth=("neo4j", "password"))

def get_friends(tx, nombre):
    query = "MATCH (p:Persona {nombre: $nombre})-[:AMIGO_DE]->(amigo) RETURN amigo.nombre"
    result = tx.run(query, nombre=nombre)
    return [record["amigo.nombre"] for record in result]

with driver.session() as session:
    amigos = session.execute_read(get_friends, "Ana")
    print(amigos)

driver.close()
    \end{lstlisting}
\end{frame}

% ============================================================================
% SECCIÓN 5: BASES DE DATOS COLUMNARES: CASSANDRA
% ============================================================================
\section{Cassandra}

\begin{frame}[fragile]{Arquitectura y Modelo de Datos}
    \begin{itemize}
        \item Peer-to-peer, sin maestro.
        \item Escalabilidad lineal.
        \item Clave primaria compuesta: \textbf{clave de partición} + \textbf{columnas de clustering}.
    \end{itemize}
    \begin{exampleblock}{Tabla para mensajes por usuario}
        \begin{lstlisting}
CREATE TABLE mensajes (
    usuario_id UUID,
    timestamp TIMESTAMP,
    mensaje TEXT,
    PRIMARY KEY ((usuario_id), timestamp)
) WITH CLUSTERING ORDER BY (timestamp DESC);
        \end{lstlisting}
    \end{exampleblock}
\end{frame}

\begin{frame}[fragile]{Conexión desde Python (cassandra-driver)}
    \begin{lstlisting}[style=python]
from cassandra.cluster import Cluster

cluster = Cluster(['127.0.0.1'])
session = cluster.connect('mi_keyspace')

rows = session.execute("SELECT * FROM mensajes WHERE usuario_id = %s", [usuario_id])
for row in rows:
    print(row.timestamp, row.mensaje)

cluster.shutdown()
    \end{lstlisting}
\end{frame}

% ============================================================================
% SECCIÓN 6: MOTORES DE BÚSQUEDA: ELASTICSEARCH
% ============================================================================
\section{Elasticsearch}

\begin{frame}{Conceptos Clave}
    \begin{itemize}
        \item \textbf{Índice}: colección de documentos.
        \item \textbf{Documento}: unidad de información (JSON).
        \item \textbf{Shard}: partición de un índice.
        \item \textbf{Inverted index}: estructura para búsqueda rápida de términos.
    \end{itemize}
\end{frame}

\begin{frame}[fragile]{Consultas con Python (elasticsearch-py)}
    \begin{lstlisting}[style=python]
from elasticsearch import Elasticsearch
from datetime import datetime

es = Elasticsearch(['localhost:9200'])

# Indexar documento
doc = {
    'timestamp': datetime.now(),
    'nivel': 'ERROR',
    'mensaje': 'Timeout en conexión'
}
es.index(index='logs', body=doc)

# Buscar
query = {
    'query': {
        'match': {'mensaje': 'Timeout'}
    }
}
resp = es.search(index='logs', body=query)
for hit in resp['hits']['hits']:
    print(hit['_source'])
    \end{lstlisting}
\end{frame}

% ============================================================================
% SECCIÓN 7: COMPARATIVA SQL vs NoSQL
% ============================================================================
\section{Comparativa SQL vs NoSQL}

\begin{frame}{Comparativa general}
    \begin{table}
        \centering
        \begin{tabular}{l|l|l}
            \toprule
            Característica & SQL & NoSQL \\
            \midrule
            Esquema & Fijo (schema-on-write) & Flexible (schema-on-read) \\
            Transacciones & ACID & BASE (consistencia eventual) \\
            Escalabilidad & Vertical & Horizontal \\
            Consultas & Estructuradas (JOINs) & Según tipo (no joins) \\
            Casos de uso & Transaccionales & Big Data, tiempo real \\
            \bottomrule
        \end{tabular}
    \end{table}
\end{frame}

\begin{frame}{Comparativa entre tipos NoSQL}
    \begin{table}
        \centering
        \begin{tabular}{l|l|l}
            \toprule
            Tipo & Ventajas & Casos típicos \\
            \midrule
            Documental & Flexible, fácil & Catálogos, CMS \\
            Clave-Valor & Ultra rápido & Caché, sesiones \\
            Grafos & Relaciones complejas & Recomendación, fraude \\
            Columnar & Alta escalabilidad & IoT, series temporales \\
            Búsqueda & Texto completo & Logs, búsqueda web \\
            \bottomrule
        \end{tabular}
    \end{table}
\end{frame}

% ============================================================================
% SECCIÓN 8: INTEGRACIÓN CON INTELIGENCIA ARTIFICIAL
% ============================================================================
\section{Integración con IA}

\begin{frame}[fragile]{Almacenamiento de Embeddings}
    \begin{itemize}
        \item En MongoDB se guardan como arrays y se usa \texttt{\$cosineSimilarity}.
        \item En Elasticsearch con campo \texttt{dense\_vector}.
        \item Ejemplo con MongoDB:
    \end{itemize}
    \begin{lstlisting}[style=python]
embedding = [0.12, 0.45, ...]  # vector de 300 dimensiones
coleccion.insert_one({'texto': 'ejemplo', 'embedding': embedding})
    \end{lstlisting}
\end{frame}

\begin{frame}[fragile]{Caché de Modelos y Resultados con Redis}
    \begin{lstlisting}[style=python]
import redis, pickle
r = redis.Redis()

# Cachear modelo serializado
modelo_bytes = pickle.dumps(modelo_entrenado)
r.set('modelo:recomendacion', modelo_bytes)

# Cachear predicción con TTL
r.setex('pred:user:123', 3600, pickle.dumps(prediccion))
    \end{lstlisting}
\end{frame}

\begin{frame}[fragile]{Sesiones de Usuario en Chatbots (Redis)}
    \begin{lstlisting}[style=python]
# Guardar contexto de conversación
import json
r.setex(f'session:{session_id}', 1800, json.dumps({
    'ultimo_intento': 'hola',
    'historial': ['mensaje1', 'mensaje2']
}))
    \end{lstlisting}
\end{frame}

\begin{frame}[fragile]{Recomendaciones con Grafos (Neo4j)}
    \begin{block}{Recomendar productos basados en usuarios similares}
        \begin{lstlisting}
MATCH (u:Usuario {id:123})-[:COMPRO]->(p:Producto)<-[:COMPRO]-(otros:Usuario)
MATCH (otros)-[:COMPRO]->(recom:Producto)
WHERE NOT (u)-[:COMPRO]->(recom)
RETURN recom, COUNT(*) AS frecuencia
ORDER BY frecuencia DESC
LIMIT 10
        \end{lstlisting}
    \end{block}
\end{frame}

% ============================================================================
% SECCIÓN 9: INFRAESTRUCTURA EN LA NUBE PARA NOSQL
% ============================================================================
\section{Infraestructura Cloud para NoSQL}

\begin{frame}{Servicios gestionados en AWS}
    \begin{description}
        \item[DynamoDB] Clave-valor/documental serverless.
        \item[DocumentDB] Compatible con MongoDB.
        \item[ElastiCache] Redis/Memcached gestionado.
        \item[Neptune] Base de datos de grafos.
        \item[OpenSearch] Sucesor de Elasticsearch.
    \end{description}
\end{frame}

\begin{frame}{Servicios gestionados en Azure}
    \begin{description}
        \item[Cosmos DB] Multimodelo (documental, clave-valor, columna, grafo).
        \item[Table Storage] Clave-valor.
        \item[Cache for Redis] Redis gestionado.
    \end{description}
\end{frame}

\begin{frame}{Servicios gestionados en GCP}
    \begin{description}
        \item[Firestore] Documental.
        \item[Bigtable] Columna ancha (similar a HBase).
        \item[Memorystore] Redis.
    \end{description}
\end{frame}

\begin{frame}{Dimensionamiento y Costos (Ejemplo DynamoDB)}
    \begin{itemize}
        \item \textbf{RCU/WCU}: Unidades de capacidad de lectura/escritura.
        \item Ejemplo: 1 millón escrituras/mes = \$1.25, 10 millones lecturas/mes = \$2.50, 100 GB almacenamiento = \$25.
        \item Total aproximado: \$28.75/mes.
    \end{itemize}
\end{frame}

% ============================================================================
% SECCIÓN 10: EJERCICIOS RESUELTOS
% ============================================================================
\section{Ejercicios Resueltos}

\begin{frame}[fragile]{Ejercicio 1: Catálogo de productos en MongoDB}
    \textbf{Enunciado:} Modelar un documento para productos con nombre, precio, categoría, etiquetas y especificaciones variables.
    \begin{block}{Solución}
        \begin{lstlisting}[style=python]
from pymongo import MongoClient

client = MongoClient()
db = client['tienda']
productos = db['productos']

producto1 = {
    'nombre': 'Laptop Gamer',
    'precio': 1200,
    'categoria': 'Electrónica',
    'etiquetas': ['laptop', 'gaming'],
    'especificaciones': {'ram': '16GB', 'procesador': 'i7'}
}
producto2 = {
    'nombre': 'Mouse',
    'precio': 25,
    'categoria': 'Electrónica',
    'etiquetas': ['mouse', 'inalámbrico'],
    'especificaciones': {'tipo': 'óptico'}
}
productos.insert_many([producto1, producto2])

result = productos.find({'categoria': 'Electrónica', 'precio': {'$gt': 100}})
for p in result:
    print(p['nombre'])
        \end{lstlisting}
    \end{block}
\end{frame}

\begin{frame}[fragile]{Ejercicio 2: Sistema de votación con Redis}
    \textbf{Enunciado:} Usar Redis para contador de votos y ranking con sorted set.
    \begin{block}{Solución}
        \begin{lstlisting}[style=python]
import redis

r = redis.Redis()

# Votar
r.incr('votos:candidato:A')
r.incr('votos:candidato:B')
r.incr('votos:candidato:A')

# Ranking
r.zincrby('ranking', 1, 'A')
r.zincrby('ranking', 2, 'B')
r.zincrby('ranking', 1, 'A')

print(r.get('votos:candidato:A'))
print(r.zrevrange('ranking', 0, -1, withscores=True))
        \end{lstlisting}
    \end{block}
\end{frame}

\begin{frame}[fragile]{Ejercicio 3: Detección de fraude con grafos (Neo4j)}
    \textbf{Enunciado:} Detectar ciclos de transferencias rápidas entre cuentas.
    \begin{block}{Solución}
        \begin{lstlisting}
-- Crear cuentas y transferencias
CREATE (a:Cuenta {id: 'A'}), (b:Cuenta {id: 'B'}), (c:Cuenta {id: 'C'})
CREATE (a)-[:TRANSFIERE {monto: 1000, fecha: '2025-03-01'}]->(b)
CREATE (b)-[:TRANSFIERE {monto: 900, fecha: '2025-03-02'}]->(c)
CREATE (c)-[:TRANSFIERE {monto: 800, fecha: '2025-03-03'}]->(a)

-- Buscar ciclos de longitud 3
MATCH p = (a)-[:TRANSFIERE*3]-(a)
RETURN p
        \end{lstlisting}
    \end{block}
\end{frame}

\begin{frame}[fragile]{Ejercicio 4: Logs en Elasticsearch}
    \textbf{Enunciado:} Indexar logs y buscar errores de los últimos 5 minutos.
    \begin{block}{Solución}
        \begin{lstlisting}[style=python]
from elasticsearch import Elasticsearch
from datetime import datetime, timedelta

es = Elasticsearch()

doc = {
    'timestamp': datetime.now(),
    'nivel': 'ERROR',
    'mensaje': 'Timeout en conexión'
}
es.index(index='logs', body=doc)

hace_5min = datetime.now() - timedelta(minutes=5)
query = {
    'query': {
        'bool': {
            'must': [
                {'match': {'nivel': 'ERROR'}},
                {'range': {'timestamp': {'gte': hace_5min}}}
            ]
        }
    }
}
resp = es.search(index='logs', body=query)
for hit in resp['hits']['hits']:
    print(hit['_source'])
        \end{lstlisting}
    \end{block}
\end{frame}

% ============================================================================
% SECCIÓN 11: GLOSARIO
% ============================================================================
\section{Glosario}

\begin{frame}{Términos Clave}
    \begin{description}
        \item[NoSQL] Not Only SQL, bases de datos no relacionales.
        \item[Teorema CAP] Consistencia, Disponibilidad, Tolerancia a Particiones.
        \item[MongoDB] Base de datos documental.
        \item[Redis] Base de datos en memoria clave-valor.
        \item[Neo4j] Base de datos de grafos.
        \item[Cassandra] Base de datos columnar distribuida.
        \item[Elasticsearch] Motor de búsqueda y analítica.
        \item[DynamoDB] NoSQL serverless de AWS.
        \item[Cosmos DB] Base de datos multimodelo de Azure.
        \item[Sharding] Particionamiento horizontal.
        \item[Replicación] Copia de datos en varios nodos.
        \item[Embedding] Representación vectorial para IA.
    \end{description}
\end{frame}

% --- Diapositiva final ---
\begin{frame}
    \centering
    \Huge \textbf{¿Preguntas?}
    
    \vspace{1cm}
    \large ¡Gracias por su atención!
\end{frame}

\end{document}

% latexmk -pdf ppt_test.tex && rm -f *.aux *.log *.out *.toc *.nav *.snm *.fdb_latexmk *.fls *.vrb